% Part: normal-modal-logic
% Chapter: filtrations
% Section: more-filtrations

\documentclass[../../../include/open-logic-section]{subfiles}

\begin{document}

\olfileid[fr]{nml}{fil}{acc}

\olsection{Filtrations et propriétés de l'accessibilité}

Comme nous l'avons vu, toute filtration d'un modèle universel, sériel
ou réflexif est, respectivement, universelle, sérielle ou réflexive.
En revanche, une filtration d'un modèle symétrique ou transitif n'est
pas nécessairement symétrique ou transitive. Dans certains cas, on
peut néanmoins définir une filtration qui conserve ces propriétés.
Pour cela, suivons la définition de la filtration la plus grossière,
en ajoutant des conditions à la définition de~$R^*$. Supposons
$\Gamma$ clos par sous-!!{formula}s. Considérons les
conditions~$C_i(u,v)$ du \olref{tab:Cn-filtrations}, portant sur des
mondes $u$, $v$ d'un modèle $\mModel{M} =\tuple{W, R, V}$. Nous
pouvons définir $R^*[u][v]$ par des combinaisons de ces conditions.
Si, par exemple, nous posons $R^*[u][v]$ si et seulement si
$C_1(u,v)$, nous retrouvons exactement la filtration la plus
grossière. Si le modèle est symétrique et si nous posons $R^*[u][v]$ si et seulement si
$C_1(u,v)$ et $C_2(u, v)$, nous obtenons une filtration « au moins aussi
fine » que la plus grossière : pas davantage de couples de mondes satisfont
conjointement $C_1(u,v)$ et $C_2(u,v)$ que la seule condition
$C_1(u,v)$. Les deux relations peuvent coïncider.

\begin{table}[ht]
  \centering
  \begin{tabular}{|ll|}
    \hline
    \multirow{2}{*}{$C_1(u,v)$:}
    \iftag{prvBox}{&
      si $\Box!A \in \Gamma$ et $\mSat{M}{\Box!A}[u]$,
      alors $\mSat{M}{!A}[v]$ ;\iftag{prvDiamond}{ et}{}\\}{}
    \iftag{prvDiamond}{&
      si $\Diamond!A \in \Gamma$ et $\mSat{M}{!A}[v]$,
      alors $\mSat{M}{\Diamond!A}[u]$ ; \\}{}
    \hline
    \multirow{2}{*}{$C_2(u,v)$:}
    \iftag{prvBox}{&
      si $\Box!A \in \Gamma$ et $\mSat{M}{\Box!A}[v]$,
      alors $\mSat{M}{!A}[u]$ ;\iftag{prvDiamond}{ et}{}\\}{}
    \iftag{prvDiamond}{&
      si $\Diamond!A \in \Gamma$ et $\mSat{M}{!A}[u]$,
      alors $\mSat{M}{\Diamond!A}[v]$ ; \\}{}
    \hline
    \multirow{2}{*}{$C_3(u,v)$:}
    \iftag{prvBox}{&
      si $\Box!A \in \Gamma$ et $\mSat{M}{\Box!A}[u]$,
      alors $\mSat{M}{\Box!A}[v]$ ;\iftag{prvDiamond}{ et}{}\\}{}
    \iftag{prvDiamond}{&
      si $\Diamond!A \in \Gamma$ et $\mSat{M}{\Diamond!A}[v]$,
      alors $\mSat{M}{\Diamond!A}[u]$ ; \\}{}
    \hline
    \multirow{2}{*}{$C_4(u,v)$:}
    \iftag{prvBox}{&
      si $\Box!A \in \Gamma$ et $\mSat{M}{\Box!A}[v]$,
      alors $\mSat{M}{\Box!A}[u]$ ;\iftag{prvDiamond}{ et}{}\\}{}
    \iftag{prvDiamond}{&
      si $\Diamond!A \in \Gamma$ et $\mSat{M}{\Diamond!A}[u]$,
      alors $\mSat{M}{\Diamond!A}[v]$ ; \\}{}
    \hline
    \end{tabular}
  \caption{Conditions entre mondes possibles pour définir des
    filtrations.}
  \ollabel{tab:Cn-filtrations}
\end{table}

\begin{thm}\ollabel{thm:more-filtrations}
  Soit $\mModel{M} =\tuple{W,R,V}$ un modèle, avec $\Gamma$ clos
  par sous-!!{formula}s. Supposons $W^*$ et $V^*$ définis comme dans la
  \olref[fil]{defn:filtration}. Alors :
  \begin{enumerate}
  \item Supposons $R^*[u][v]$ si et seulement si $C_1(u, v) \land
    C_2(u,v)$. Alors $R^*$ est symétrique, et $\mModel{M^*} =
    \tuple{W^*,R^*,V^*}$ est une filtration si $\mModel{M}$ est
    symétrique.
  \item Supposons $R^*[u][v]$ si et seulement si $C_1(u, v)
    \land C_3(u,v)$. Alors $R^*$ est transitive, et
    $\mModel{M^*}=\tuple{W^*,R^*,V^*}$ est une filtration si
    $\mModel{M}$ est transitif.
  \item Supposons $R^*[u][v]$ si et seulement si $C_1(u, v) \land C_2(u,v)
    \land C_3(u,v) \land C_4(u,v)$. Alors $R^*$ est symétrique et
    transitive, et $\mModel{M^*}=\tuple{W^*,R^*,V^*}$ est une
    filtration si $\mModel{M}$ est symétrique et transitif.
  \item Supposons que $R^*$ soit définie par $R^*[u][v]$ si et
    seulement si $C_1(u,
    v) \land C_3(u,v) \land C_4(u,v)$. Alors $R^*$ est transitive et
    euclidienne, et $\mModel{M^*}=\tuple{W^*,R^*,V^*}$ est une
    filtration si $\mModel{M}$ est transitif et euclidien.
  \end{enumerate}
\end{thm}

\begin{proof}
  \begin{enumerate}
    \item La symétrie de $R^*$ est immédiate, puisque $C_1(u,v)
      \Leftrightarrow C_2(v,u)$ et $C_2(u,v) \Leftrightarrow
      C_1(v,u)$. Il reste donc à montrer que, si $\mModel{M}$ est
      symétrique, $\mModel{M^*}$ est une filtration par~$\Gamma$.
      La condition $C_1(u,v)$ garantit que
      \iftag{prvBox}{%
        \iftag{prvDiamond}{%
          \olref[fil]{defn:filtration-R2} et
          \olref[fil]{defn:filtration-R3} de la
          \olref[fil]{defn:filtration} sont satisfaites}{%
          la \olref[fil]{defn:filtration}\olref[fil]{defn:filtration-R2} est satisfaite}}{%
        la \olref[fil]{defn:filtration}\olref[fil]{defn:filtration-R3} est satisfaite}.
      Il suffit donc de vérifier la
      \olref[fil]{defn:filtration}\olref[fil]{defn:filtration-R1},
      c'est-à-dire que $Ruv$ entraîne $R^*[u][v]$.

      Supposons donc $Ruv$. Pour établir $R^*[u][v]$, il faut
      vérifier $C_1(u,v)$ et $C_2(u,v)$. Pour $C_1$ :
      \iftag{prvBox}{si $\Box!A
        \in\Gamma$ et $\mSat{M}{\Box!A}[u]$, alors
        $\mSat{M}{!A}[v]$ puisque $Ruv$.\iftag{prvDiamond}{ De même,
        }{}}{}\iftag{prvDiamond}{si $\Diamond!A \in \Gamma$ et
        $\mSat{M}{!A}[v]$, alors $\mSat{M}{\Diamond!A}[u]$ puisque
        $Ruv$.}{} Pour $C_2$ : \iftag{prvBox}{si $\Box!A \in \Gamma$ et
        $\mSat{M}{\Box!A}[v]$, alors $Ruv$ entraîne $Rvu$ par symétrie,
        d'où $\mSat{M}{!A}[u]$.\iftag{prvDiamond}{ De même,
        }{}}{}\iftag{prvDiamond}{si $\Diamond!A \in\Gamma$ et
        $\mSat{M}{!A}[u]$, alors $\mSat{M}{\Diamond!A}[v]$, puisque
        $Rvu$ par symétrie.}{}
    \item Exercice.
    \item Exercice.
    \item Exercice.
  \end{enumerate}
\end{proof}

\begin{prob}
  Complétez la démonstration du \olref[nml][fil][acc]{thm:more-filtrations}.
\end{prob}
\begin{explain}
Les conditions $C_1(u,v)$ et $C_2(u,v)$ s'échangent
lorsqu'on échange $u$ et~$v$. Leur conjonction définit donc
une relation symétrique. Pour obtenir une filtration, il
faut en outre conserver chaque flèche héritée de~$R$.
L'hypothèse de symétrie de~$R$ garantit ici cette conservation.
\end{explain}
\end{document}
